\section{Modulos y responsabilidades}

\subsection{Descripcion}

Este documento define los modulos esperados del backend y su reparto de responsabilidades. Su objetivo es dejar claro que le pertenece a cada modulo, que tipo de entradas y salidas maneja, y que dependencias son aceptables dentro del modular monolith.

\subsection{Alcance}

Aplica al repositorio backend y a su organizacion por dominio. No redefine reglas del negocio ya cerradas; solo distribuye ownership tecnico-funcional.

\subsection{Definiciones}

\begin{itemize}
  \item \textbf{Modulo}: unidad de organizacion del backend alrededor de una responsabilidad de dominio.
  \item \textbf{Ownership}: responsabilidad principal de un modulo sobre una parte del comportamiento o de los datos del sistema.
  \item \textbf{Dependencia permitida}: relacion entre modulos aceptable dentro del diseno del backend.
\end{itemize}

\subsection{Reglas}

\begin{itemize}
  \item El backend se documentara por modulos de dominio.
  \item Cada modulo debera declarar sus entradas, salidas y dependencias permitidas.
  \item La logica de negocio no se documentara como responsabilidad de controllers.
  \item Un modulo no debe absorber responsabilidades ajenas solo por conveniencia de implementacion.
\end{itemize}

\subsubsection{Tree modular esperado}

\begin{verbatim}
modules/
|-- auth/
|-- users/
|-- payments/
|-- matches/
|-- picks/
|-- results/
|-- scoring/
|-- leaderboard/
|-- admin/
|-- audit/
`-- notifications/
\end{verbatim}

\subsubsection{Responsabilidades por modulo}

\begin{longtable}{p{0.16\textwidth}p{0.24\textwidth}p{0.23\textwidth}p{0.25\textwidth}}
\toprule
\textbf{Modulo} & \textbf{Ownership principal} & \textbf{Entradas / salidas esperadas} & \textbf{Dependencias permitidas} \\
\midrule
\texttt{auth} & autenticacion, sesion y acceso inicial & credenciales, tokens, sesion valida o rechazada & \texttt{users}, servicios transversales \\
\texttt{users} & identidad y estado funcional de cuenta & registro, consulta y cambio de estado de usuario & \texttt{audit}, \texttt{notifications} \\
\texttt{payments} & comprobantes, revision y activacion administrativa & carga y revision de comprobante, resolucion administrativa & \texttt{users}, \texttt{audit}, \texttt{notifications} \\
\texttt{matches} & calendario, fases, kickoff y estado de partido & partidos, fases, horarios y cambios de estado & \texttt{audit} \\
\texttt{picks} & captura, edicion y lock de picks y predicciones especiales & picks validos o rechazados, estado de lock & \texttt{users}, \texttt{matches}, \texttt{audit} \\
\texttt{results} & resultado vigente, sincronizacion externa y overrides & resultados externos o manuales, verdad oficial vigente & \texttt{matches}, \texttt{audit} \\
\texttt{scoring} & calculo de puntos, multiplicadores y desempates & score entries, recalculo y total acumulado & \texttt{picks}, \texttt{results}, \texttt{matches}, \texttt{audit} \\
\texttt{leaderboard} & vistas derivadas del ranking & ranking general, por jornada o fase & \texttt{scoring}, \texttt{users} \\
\texttt{admin} & coordinacion de acciones administrativas autorizadas & comandos administrativos, ejecucion auditada & cualquier modulo de dominio por casos de uso explicitos \\
\texttt{audit} & trazabilidad de acciones y procesos & eventos auditables persistidos y consultables & servicio transversal observado por otros modulos \\
\texttt{notifications} & envio de correos y avisos funcionales & mensajes y disparadores de comunicacion & \texttt{users}, \texttt{payments}, \texttt{results}, \texttt{admin} \\
\bottomrule
\end{longtable}

\subsubsection{Reglas de colaboracion entre modulos}

\begin{itemize}
  \item \texttt{payments} no activa usuarios por mutacion directa informal; la activacion debe quedar trazada y reflejada en \texttt{users}.
  \item \texttt{picks} controla captura y lock, pero no calcula puntaje.
  \item \texttt{results} decide la verdad vigente del resultado, pero no recalcula ranking por si solo; esa responsabilidad pertenece a \texttt{scoring} y \texttt{leaderboard}.
  \item \texttt{leaderboard} consume resultados de scoring y no redefine desempates ni formulas de puntos.
  \item \texttt{admin} coordina, no reemplaza ownership de los modulos de dominio.
\end{itemize}

\subsection{Casos borde}

\begin{itemize}
  \item Si una accion administrativa impacta varios modulos, cada cambio relevante debe seguir perteneciendo a su modulo duenio y quedar auditado.
  \item Si \texttt{notifications} se implementa como servicio transversal en lugar de modulo independiente, su ownership funcional debe seguir documentado para no quedar difuso.
\end{itemize}

\subsection{Invariantes}

\begin{itemize}
  \item Ningun modulo debe redefinir reglas centrales de \texttt{01\_Producto\_y\_Reglas}.
  \item Ningun controller debe convertirse en duenio efectivo de reglas de dominio.
  \item Toda colaboracion entre modulos debe preservar un ownership principal identificable.
\end{itemize}

\subsection{Preguntas abiertas}

\begin{itemize}
  \item Confirmar la frontera final entre \texttt{admin}, \texttt{audit} y \texttt{notifications} cuando se documenten procesos backend detallados.
\end{itemize}

\subsection{Decisiones pendientes}

\begin{itemize}
  \item `Decision pendiente` `No bloqueante`: decidir si \texttt{notifications} queda como modulo propio o como capacidad transversal con ownership distribuido.
  \item `Decision pendiente` `No bloqueante`: confirmar si \texttt{audit} expone consultas administrativas como modulo propio o solo como infraestructura observable por otros modulos.
\end{itemize}
